% !TEX root = approx8.tex

\appendix

\section{Filtered $A_{\infty}$-categories} \label{a:fil-ai}
\subsection{Filtered $A_\infty$-categories, $A_\infty$-functors, $A_\infty$-natural transformations and shift $A_\infty$-functors}\label{a:sb:fil-ai}
 Filtered
$A_{\infty}$-categories have already appeared in the literature, see
e.g.~\cite[Section~3.2]{BCZ:tpc}
and~\cite[Chapter~2]{Bi-Co-Sh:LagrSh}. We will therefore go over this
subject very briefly and cover below mainly aspects of the theory that
have not been addressed in the literature.

Fix a field $\k$ and $\k_0 \subset \k$ a subring. For most of the
applications in this paper we will either have
$\k_0 = \k = \mathbb{Z}_2$, or $\k = \Lambda$ (the universal Novikov
field with coefficients in $\mathbb{Z}_2$) and
$\k_0 = \Lambda_0 \subset \Lambda$ the positive Novikov ring
(see~\eqref{eq:Nov-ring} and~\eqref{eq:Nov-ring-0}
in~\S\ref{sb:fuk-mon}).

In the followings, we say that $(C,d)$ is a filtered chain complex over $\k$ if $(C,d)$ is chain complex over $\k$ and comes endowed with an increasing filtration parametrized by the real line of $\k_0$-modules that are preserved by $d$.

Let $\mathcal{A}$ be an $A_{\infty}$-category over $\k$. For
$X_0, X_1 \in \Ob(\mathcal{A})$ we abbreviate:
$$\mathcal{A}(X_0,X_1):= \hom_{\mathcal{A}}(X_0,X_1),$$ and
more generally, for a tuple of objects
$\vec{X} := (X_0, \ldots, X_d)$, $d \geq 1$, we write
$$\mathcal{A}(\vec{X}) := \mathcal{A}(X_0, X_1) \otimes_{\k} \cdots
\otimes_{\k} \mathcal{A}(X_{d-1}, X_d)$$ and denote by
$\mu_d: \mathcal{A}(\vec{X}) \longrightarrow \mathcal{A}(X_0, X_d)$
the $A_{\infty}$-operations.

A filtered $A_{\infty}$-category (over $(\k,\k_0)$) is an
$A_{\infty}$-category $\mathcal{A}$ over $\k$ such that the spaces of
morphisms $\mathcal{A}(X, Y) = \hom_{\mathcal{A}}(X, Y)$ between every
two objects $X, Y$ are filtered chain complexes over $\k$
%  are filtered by an increasing filtration of $\k_0$-modules parametrized by $\mathbb{R}$, and such that all the
%$A_{\infty}$-operations $\mu_d$, $d \geq 1$, preserve the filtration
%levels.}The precise definition is as follows. 
For every
$\alpha \in \mathbb{R}$ we denote by
$\mathcal{A}^{\alpha}(X,Y) \subset \mathcal{A}(X,Y)$ the
$\k_0$-submodule which is the filtration level parametrized by
$\alpha$. More generally, for $\alpha \in \mathbb{R}$ and a tuple
$\vec{X} = (X_0, \ldots, X_d)$ of objects of $\mathcal{A}$ we write
$$\mathcal{A}^\alpha(\vec{X}) \subset \mathcal{A}(\vec{X})$$ for the
$\k_0$-submodule generated by the images of all the $\k_0$-modules
$$\mathcal{A}^{\alpha_1}(X_0, X_1) \otimes_{\k_0} \cdots \otimes_{\k_0}
\mathcal{A}^{\alpha_d}(X_{d-1}, X_d), \text{ with } \alpha_1 + \cdots
+ \alpha_d \leq \alpha,$$ under the obvious ($\k_0$-linear) map
$$\mathcal{A}(X_0, X_1) \otimes_{\k_0} \cdots \otimes_{\k_0}
\mathcal{A}(X_{d-1}, X_d) \longrightarrow \mathcal{A}(\vec{X}).$$ We
then require that for every $d \geq 1$, every $\vec{X}$ and every
$\alpha \in \mathbb{R}$, we have:
$$\mu_d (\mathcal{A}^{\alpha}(\vec{X}))
\subset \mathcal{A}^{\alpha}(X_0, X_d).$$

% and for every tuple of objects
% $\vec{X} = (X_0, \ldots, X_d)$ of $\mathcal{A}$ we write
% $\mathcal{A}^{\alpha}(\vec{X})$ for the $\k_0$-submodule of
% $\mathcal{A}(\vec{X})$ generated by all the submodules
% $$\mathcal{A}^{\alpha_1}(X_0, X_1) \otimes_{\k_0} \cdots \otimes_{\k_0}
% \mathcal{A}^{\alpha_d}(X_{d-1}, X_d) \subset \mathcal{A}(\vec{X}), \;
% \text{with} \; \alpha_1 + \cdots + \alpha_d \leq \alpha.$$ Then we
% require that for every $d \geq 1$ and every $\vec{X}$ we have
% $$\mu_d (\mathcal{A}^{\alpha}(\vec{X}))
% \subset \mathcal{A}^{\alpha}(X_0, X_d), \; \forall \alpha \in
% \mathbb{R}.$$ Another useful piece of notation is the following.
% Consider a tuple $\vec{\alpha} = (\alpha_1, \ldots, \alpha_d)$,
% $d \geq 1$, $\alpha_i \in \mathbb{R}$ of real numbers and denote
% $|\vec{\alpha}| = \alpha_1 + \cdots + \alpha_d$. For a tuple $\vec{X}$
% of objects in $\mathcal{A}$ define
% $$\mathcal{A}^{\vec{\alpha}}(\vec{X}) :=
% \mathcal{A}^{\alpha_1}(X_0, X_1) \otimes_{\k_0} \cdots \otimes_{\k_0}
% \mathcal{A}^{\alpha_d}(X_{d-1}, X_d) \subset \mathcal{A}(\vec{X}).$$

A filtered $A_{\infty}$-category $\mathcal{A}$ is called strictly
unital (in the filtered sense) if it is strictly unital as an
$A_{\infty}$-category without filtrations and moreover for every
$X \in \Ob(\mathcal{A})$ the strict unit $e_X \in \mathcal{A}(X,X)$
lies in filtration level $0$, i.e.~$e_X \in \mathcal{A}^0(X,X)$.
Unless otherwise stated, all the filtered $A_{\infty}$-categories
$\mathcal{A}$ in this paper are implicitly assumed to be strictly
unital.

An example of filtered $A_\infty$-category is the $dg$-category $\mathcal{F}\mathbf{Ch}$ of filtered chain complexes over $\k$.

The concept of $A_{\infty}$-functors has a filtered counterpart.  An
$A_{\infty}$-functor
$\mathcal{F}: \mathcal{A} \longrightarrow \mathcal{B}$ between two
$A_{\infty}$-categories is said to be filtered if for every $d \geq 1$
and every tuple of objects $\vec{X}$ its $d$-th order term
$\mathcal{F}_d$ preserves filtrations in the sense that
$$\mathcal{F}_d(\mathcal{A}^{\alpha}(\vec{X})) \subset
\mathcal{B}^{\alpha}(\mathcal{F}X_0, \mathcal{F}X_d), \; \forall
\alpha \in \mathbb{R}.$$ All the filtered $A_{\infty}$-functors in
this paper will be implicitly assumed to be strictly unital.

One can define pre-natural transformations between filtered
$A_{\infty}$-functors as follows.  Let
$\mathcal{F}, \mathcal{G}: \mathcal{A} \longrightarrow \mathcal{B}$ be
two filtered $A_{\infty}$-functors. A pre-natural transformation
$T: \mathcal{F} \longrightarrow \mathcal{G}$ of shift $s$ is a
pre-natural transformation (in the usual $A_{\infty}$-sense) such that
its $d$-th order term $T_d$ shifts filtrations by not more than $s$, namely we have
a family of elements
$T^0 \in \mathcal{B}^{s}(\mathcal{F} X, \mathcal{G}X)$ for every
$X \in \Ob(\mathcal{A})$ and moreover, for every $d \geq 0$ and every
tuple of objects $\vec{X} = (X_0, \ldots, X_d)$ we have:
$$T_d (\mathcal{A}^{\alpha}(\vec{X})) \subset
\mathcal{B}^{\alpha + s}(\mathcal{F}X_0, \mathcal{G}X_d), \; \forall d
\geq 1, \, \alpha \in \mathbb{R}.$$

We will always endow our $A_{\infty}$-categories $\mathcal{A}$ with a
shift functor. This is not a single functor but in fact a system of
filtered $A_{\infty}$-functors
$\{\Sigma^r : \mathcal{A} \longrightarrow \mathcal{A}\}_{r \in
  \mathbb{R}}$, $r \in \mathbb{R}$, all having trivial higher order
terms ($\Sigma^r_d = 0, \, \forall d \geq 2$), and together with
filtered natural transformations
$\eta_{r,s}: \Sigma^r \longrightarrow \Sigma^s$ of shift $s-r$, for
every $r,s \in \mathbb{R}$, and such that:
\begin{enumerate}
\item $\Sigma^0 = \id$. \label{i:shift-a}
\item $\Sigma^r \circ \Sigma^s = \Sigma^{r+s}$ for all $r,s \in \mathbb{R}$.
\item The higher order terms $(\eta_{r,s})_d, d \geq 1$ of
  $\eta_{r,s}$ vanish.
\item For every $X \in \Ob(\mathcal{A})$ the $0$-term of $\eta_{r,s}$
  assigned to $X$ is a cycle
  $(\eta^X_{r,s})_0 \in \mathcal{A}^{s-r}(\Sigma^rX, \Sigma^sX)$.
  Moreover, we have $(\eta^X_{0,0})_0 = e_X \in \mathcal{A}^0(X,X)$.
\item For every $X \in \Ob(\mathcal{A})$ and every $t \in \mathbb{R}$
  we have
  $\mu_2 \bigl( (\eta^X_{r,s})_0, (\eta^X_{s,t})_0 \bigr) =
  (\eta^X_{r, t})_0$.
\item For every $X \in \Ob(\mathcal{A})$ and every $t \in \mathbb{R}$
  we have
  $\Sigma_1^t \bigl( (\eta^X_{r,s})_0 \bigr) = (\eta^X_{r+t, s+t})_0$.
  \label{i:shift-b}
\end{enumerate}
It follows that each $\eta_{r,s}$ is an isomorphism and its inverse is
$\eta_{s,r}$. By composing from the left and from the right with
$\eta^X_{0,r}$ and $\eta^Y_{s,0}$ respectively we obtain isomorphisms
$$\mathcal{A}^{\alpha}(\Sigma^rX, \Sigma^sY) \cong
\mathcal{A}^{\alpha+r-s}(X,Y), \; \forall \, \alpha \in \mathbb{R}$$
that are compatible with the filtration
$\mathcal{A}^{\beta'}(- , -) \subset \mathcal{A}^{\beta''}(-,-)$,
$\beta' \leq \beta''$.


\begin{rem}
  The conditions~\eqref{i:shift-a}-\eqref{i:shift-b} may look
  excessively strong in the $A_{\infty}$-context and can probably be
  relaxed to give a weaker yet meaningful definition of shift
  functors. A more appropriate name for $\Sigma$ as above might be a
  {\em strict} shift functor, but we will not use this term in this
  paper.
\end{rem}

From now on all filtered $A_{\infty}$-functors will be assumed to be
compatible with the shift functor in the sense that
$\mathcal{F} \circ \Sigma^r = \Sigma^r \circ \mathcal{F}$ for every
$r \in \mathbb{R}$ and
$(\eta^{\mathcal{F} X}_{r,s})_0 = \mathcal{F}_1(\eta^{X}_{r,s})_0$ for
every $r, s \in \mathbb{R}$.


\subsection{Filtered $A_{\infty}$-modules and
  bimodules} \label{a:fmod} Let $\mathcal{A}$ be a filtered
$A_{\infty}$-category. A filtered left $\mathcal{A}$-module
$\mathcal{M}$ is an $\mathcal{A}$-module with the following additional
structures and properties. For every $X \in \Ob(\mathcal{A})$ the
chain complex $(\mathcal{M}(X), \mu^{\mathcal{M}}_1)$ (of vector
spaces over $\k$) is filtered by an increasing filtration of
subcomplexes of $\k_0$-modules which is parametrized by
$\mathbb{R}$. For $\alpha \in \mathbb{R}$ we denote by
$\mathcal{M}^{\alpha}(X) \subset \mathcal{M}(X)$ the $\alpha$-level of
this filtration. For $\alpha \in \mathbb{R}$ and a tuple of objects
$\vec{X} = (X_0, \ldots, X_k)$ denote by
\begin{equation} \label{eq:fil-AM}
  (\mathcal{A}(\vec{X}) \otimes \mathcal{M}(X_k))^{\alpha} \subset
  \mathcal{A}(\vec{X}) \otimes_{\k} \mathcal{M}(X_k)
\end{equation}
the $\k_0$-submodule generated by the images of all the
$\k_0$-submodules
$$\mathcal{A}^{\alpha'_1}(X_0, X_1) \otimes_{\k_0} \cdots
\otimes_{\k_0}\mathcal{A}^{\alpha'_k}(X_{k-1}, X_k) \otimes_{\k_0}
\mathcal{M}^{\alpha''}(X_k), \text{ with } \alpha'_1 + \cdots+
\alpha'_k + \alpha'' \leq \alpha,$$ under the obvious map
$$\mathcal{A}(X_0, X_1) \otimes_{\k_0} \cdots
\otimes_{\k_0}\mathcal{A}(X_{k-1}, X_k) \otimes_{\k_0}
\mathcal{M}(X_k) \longrightarrow \mathcal{A}(\vec{X}) \otimes_{\k}
\mathcal{M}(X_k).$$ We require that the higher operations
$\mu^{\mathcal{M}}_{k|1}$, $k \geq 1$, preserves filtrations. Namely, for every tuple of objects
$\vec{X} = (X_0, \ldots, X_k)$ in $\mathcal{A}$ and every
$\alpha \in \mathbb{R}$ we have:
$$\mu^{\mathcal{M}}_{k|1} \bigl( (\mathcal{A}(\vec{X}) \otimes
\mathcal{M}(X_k))^{\alpha}) \subset \mathcal{M}^{\alpha}(X_0).$$ For
uniformity we will denote the differential $\mu_1^{\mathcal{M}}$ also
by $\mu_{0|1}^{\mathcal{M}}$.\\
Alternatively, one can view filtered $\mathcal{A}$-modules as filtered $A_\infty$-functors $\mathcal{A}\to \F\textbf{Ch}^{\text{op}}$.

Pre-homomorphisms between filtered modules are defined as follows. Let
$\mathcal{M}$, $\mathcal{N}$ be filtered $A_{\infty}$-modules. A
module pre-homomorphism $t: \mathcal{M} \longrightarrow \mathcal{N}$
of shift $\leq s \in \mathbb{R}$ is a pre-homomorphism such that for
every $k \geq 0$, every tuple $\vec{X}$, and $\alpha \in \mathbb{R}$ 
we have
$$t_{k|1}( (\mathcal{A}(\vec{X}) \otimes
\mathcal{M}(X_k))^{\alpha}) \subset \mathcal{N}^{\alpha + s}(X_0).$$

Filtered $A_{\infty}$-modules and their pre-homomorphisms form a
filtered $dg$-category $F\md_{\mathcal{A}}$, where
$\hom^{s}_{F\md_{\mathcal{A}}}(\mathcal{M}, \mathcal{N})$ consists of
all pre-homomorphisms $t: \mathcal{M} \longrightarrow \mathcal{N}$ of
shift $\leq s$. Note that $F\md_{\mathcal{A}}$ comes with a natural shift functor $\Sigma$.
$\Sigma$. Its action on objects
$\mathcal{M} \in \Ob(F\md_{\mathcal{A}})$ is given by:
$$(\Sigma^r \mathcal{M})^{\alpha}(X) := \mathcal{M}^{\alpha-r}(X),
\; \forall r \in \mathbb{R}, \alpha \in \mathbb{R}, X \in
\Ob(\mathcal{A}).$$

Unless otherwise stated, all the modules in this paper will be assumed
to be strictly unital. We will also assume all the filtered modules to
be compatible with the shift functor in the obvious sense, i.e. $\mathcal{M}(\Sigma^rX) = (\Sigma^r\mathcal{M})(X)$ for all objects $X\in \Ob(\mathcal{a})$ and all $r\in \mathbb{R}$.

Right $\mathcal{A}$-modules can be defined in a completely analogous
way to left modules. To distinguish the two, whenever confusion may
arise we will denote the category of filtered left $\mathcal{A}$-modules by
$F\md^l_{\mathcal{A}}$ and the right modules by
$F\md^r_{\mathcal{A}}$.

Given two filtered $A_{\infty}$-categories $\mathcal{A}, \mathcal{B}$
we also have the filtered $(\mathcal{A}, \mathcal{B})$ modules. They
are defined by similar means to the above. Filtered bimodules too will
be assumed in this paper to be strictly unital and compatible with the
shift functor. They form a filtered $dg$-category which we
denote by $F\text{bimod}_{\mathcal{A}, \mathcal{B}}$.


\subsection{The filtered Yoneda embedding}
\subsubsection{Definition}\label{gio:fyon-def}
Let $\mathcal{A}$ be a filtered $A_\infty$-category with structure maps $(\mu_d)_{d\geq 1}$. Recall that we always assume that our $A_\infty$-categories are strictly unital. We upgrade the Yoneda embedding (see \cite[Section (1l)]{Se:book-fukaya-categ}) to a filtered version, that is, we define a filtered $A_\infty$-functor $$\mathcal{Y}\colon \mathcal{A}\to F\md_\mathcal{A}.$$ We will work with left $A_\infty$-modules, but the same can be analogously developed for right ones. Whenever confusion may arise we will denote by $\mathcal{Y}^l\colon \mathcal{A}\to F\md^l_\mathcal{A}$ the left Yoneda embedding and by $\mathcal{Y}^r\colon \mathcal{A}\to F\md_\mathcal{A}^r$. Sometimes we will need to specicify the $A_\infty$-category in the notation and we will write the Yoneda embedding as $\mathcal{Y}_\mathcal{A}$.\\

Let $L$ be an object of $\mathcal{A}$. We define $\mathcal{Y}(L)$ via $$\mathcal{Y}(L)(L'):= \mathcal{A}(L',L)$$ for any object $L'$ of $\mathcal{A}$ and endow it with the $A_\infty$-structure given by $\mu^{\mathcal{Y}(L)}_{l|1}:=\mu_{l+1}$ for $l\geq 0$. It is straightforward to see that $\mathcal{Y}(L)$ indeed defines a filtered $A_\infty$-module.\\
We now define higher operations of the functor $\mathcal{Y}$. To ease the notation, we will often write $\mathcal{Y}_L$ instead of $\mathcal{Y}(L)$ in the rest of this subsection. Consider two objects $L_0$ and $L_1$ of $\mathcal{A}$ and an element $c\in \mathcal{A}^{\leq \alpha}(L_0,L_1)$ for some $\alpha \in \mathbb{R}$. We define its image $\mathcal{Y}_1(c)\in F\md_\mathcal{A}(\mathcal{Y}_{L_0},\mathcal{Y}_{L_1})$ as the pre-morphism of modules with components $\mathcal{Y}_1(c)_{l|1}\colon \mathcal{A}(\vec{X})\otimes \mathcal{Y}_{L_0}(X_l)\to \mathcal{Y}_{L_1}(X_0)$ given by $$\mathcal{Y}_1(c)_{l|1}(x_1,\ldots, x_l,y):= \mu_{l+1}(x_1,\ldots,x_l, y,c).$$ We define the higher components $\mathcal{Y}_d$ of the Yoneda embedding in a similar manner using contraction maps, that is, given a tuple $\vec{L}=(L_0,\ldots, L_d)$ and an element $\vec{c}\in \mathcal{A}(\vec{L})$, we define the pre-morphism of modules $\mathcal{Y}_d(\vec{c})\in F\md_\mathcal{A}(\mathcal{Y}_{L_0},\mathcal{Y}_{L_d})$ via $$\mathcal{Y}_d(\vec{c})_{l|1}(x_1,\ldots, x_l,y) := \mu_{l+d+1}(x_1,\ldots, x_l,y,\vec{c}).$$ Notice that $\mathcal{Y}$ is a filtered $A_\infty$-functor, since the maps $\mu_d$, $d\geq 1$, are filtered by assumption.

%=======================
%=======================
%=======================


\subsubsection{The $\lambda$-map}\label{lambdamap}
Consider an object $L$ of $\mathcal{A}$ and a filtered $A_\infty$-module $\mathcal{M}$. We have the following map, usually simply called the $\lambda$-map: $$\lambda\colon \mathcal{M}(L)\to F\md_\mathcal{A}(\mathcal{Y}(L),\mathcal{M})$$ defined for $c\in \mathcal{M}(L)$ via $$\lambda(c)_{l|1}(x_1,\ldots, x_l,y):= \mu^\mathcal{M}_{l+1|1}(x_1,\ldots, x_l,y,c).$$ It is clear that $\lambda$ is well-defined (i.e. that it lands in the category of \textit{filtered} modules), simply because $\mathcal{A}$ is a filtered $A_\infty$-category. Similarly to \cite[Lemma 2.12]{Se:book-fukaya-categ} and \cite[Proposition 2.5.1]{Bi-Co-Sh:LagrSh} we have the following result. This result is usually proved using spectral sequence, but strict unitality of $\mathcal{A}$ allows us for a more explicit proof, which we present here.
\begin{prop}
	The map $\lambda$ is a filtered quasi-isomorphism, with filtered quasi-inverse via filtered chain-homotopies. In particular, it induces an isomorphism in persistence homology.
\end{prop}
\begin{proof}
	The fact that $\lambda$ is filtered is obvious, as $\mathcal{A}$ is filtered. We will prove the rest of the statement by explicitly constructing a quasi-inverse and the chain homotopies. First, we define a candidate quasi-inverse $$\theta\colon F\md_\mathcal{A}(\mathcal{Y}(L),\mathcal{M})\to \mathcal{M}(L)$$ as follows: given a pre-module morphism $\varphi\in F\md_\mathcal{A}(\mathcal{Y}(L),\mathcal{M}) $, we set $$\theta(\varphi):= \varphi_{0|1}(e_L)$$ where we recall that $e_L\in \mathcal{A}(L,L)$ stands for the (given) choice of strict unit for $L$. Clearly $\theta$ is a chain map. It is easy to see that given $c\in \mathcal{M}(L)$, $\theta(\lambda(c))=c$, exactly because $e_L$ is a \textit{strict} unit, so no chain homotopy is needed in this direction. Moreover, since $e_L\in \mathcal{A}^{\leq 0}(L,L)$ by our definition of filtered $A_\infty$-category, we get that $\theta$ is a filtered chain map, i.e. it sends pre-morphisms with shift $\alpha$ to $\mathcal{M}^{\leq \alpha}(L)$ for any $\alpha\in \mathbb R$.\\ We now define a candidate chain homotopy for the other composition: we define $$H\colon F\md_\mathcal{A}(\mathcal{Y}(L),\mathcal{M})\to F\md_\mathcal{A}(\mathcal{Y}(L),\mathcal{M})$$ by sending $\varphi \in F\md_\mathcal{A}(\mathcal{Y}(L),\mathcal{M})$ to the pre-morphism of modules $H(\varphi)$ with components $$H(\varphi)_{l|1}(x_1,\ldots, x_l,y):= \varphi_{l+1|1}(x_1,\ldots, x_l,y,e_L)$$ for any $l\geq 0$. Again, since strict units lie at vanishing filtration level, the map $H$ preserves $F\md_\mathcal{A}(\mathcal{Y}(L),\mathcal{M})^{\leq \alpha}$ for any $\alpha\in \mathbb R$. We prove that $$\lambda\circ \theta + id_{F\md_\mathcal{A}(\mathcal{Y}(L),\mathcal{M})} = \mu_1^{\md}\circ H + H \circ \mu_1^{\md},$$ that is, that $\theta$ is a right quasi-inverse of $\lambda$ via the chain-homotopy $H$. Let $\varphi\in F\md_\mathcal{A}(\mathcal{Y}(L),\mathcal{M})$ and $l\geq 0$, we compute:
	\begin{enumerate}
		\item first, the term \begin{align*}
			\mu_1^{\md}(H(\varphi))_{l|1}(x_1,\ldots, x_l,y)  =&  \sum_{i=0}^l \mu^\mathcal{M}_{i|1}(x_1,\ldots, x_i, \varphi_{l-i+1|1}(x_{i+1},\ldots, x_l,y,e_L)) \\ & + \sum_{i=0}^l \varphi_{i+1|1}(x_1,\ldots, x_i, \mu_{l-i+1}(x_{i+1},\ldots, x_l,y),e_L) \\
			& + \sum_{j=0}^{l-1}\sum_{i=1}^{l-j} \varphi_{l-i+2|1}(x_1,\ldots, x_j, \mu_{i}(x_{j+1},\ldots, x_{j+i}),x_{j+i+1},\ldots, x_l,y,e_L)
		\end{align*}
		\item then the term \begin{align*}
			H(\mu_1^{\md}\varphi)_{l|1}(x_1,\ldots, x_l,y)  =& \ \   \mu_1^{\md}\varphi_{l+1|1}(x_1,\ldots, x_l,y,e_L) \\  =&   \sum_{i=0}^l \mu^\mathcal{M}_{i|1}(x_1,\ldots, x_i, \varphi_{l-i+1|1}(x_{i+1},\ldots, x_l,y,e_L) + \mu_{l+1|1}^\mathcal{M}(x_1,\ldots, x_l,y,\varphi_1(e_L))\\ & +\sum_{i=0}^l \varphi_{i|1}(x_1,\ldots, x_i, \mu_{l-i+2}(x_{i+1},\ldots, x_l,y,e_L))
			\\ & + \sum_{j=0}^{l-1}\sum_{i=1}^{l-j} \varphi_{l-i+2|1}(x_1,\ldots, x_j, \mu_{}(x_{j+1},\ldots, x_{j+i}),x_{j+i+1},\ldots, x_l,y,e_L) \\ & + \sum_{j=0}^l \varphi_{j+1|1}(x_1,\ldots, x_j, \mu_{l-j+1}(x_{j+1},\ldots, x_l,y),e_L)
		\end{align*}
		and note that the second big sum after the last equality equals $\varphi_{l|1}(x_1,\ldots, x_l,y)$ since $e_L$ is a strict unit.
		\item finally, we compute \begin{align}
			(\lambda\circ \theta + id)(\varphi)_{l|1}(x_1,\ldots, x_l,y) = \mu_{l+1|1}^\mathcal{M}(x_1,\ldots, x_l,y,\varphi_1(e_L)) + \varphi_{l|1}(x_1,\ldots, x_l,y).
		\end{align}
		
	\end{enumerate}
	We see that summing all the contributions above we get the expected equality at the $l|1$ level. This proves the proposition.
\end{proof}
The following result is an easy corollary of the above. We denote by $\mathcal{Y}(\mathcal{A})=\text{image}(\mathcal{Y})$ the subcategory of $F\md_\mathcal{A}$ consisting of Yoneda modules.
\begin{cor}
	The induced persistence functor $H(\mathcal{Y})\colon H(\mathcal{A})\to H(\text{image}(\mathcal{Y}))$ is an equivalence of persistence categories.
\end{cor}



\subsection{Persistence Hochschild homology}\label{app:PHH}
Let $\mathcal{A}$ be an $A_{\infty}$-category and $\mathcal{M}$ an
$(\mathcal{A}, \mathcal{A})$-bimodule. Denote by
$CC(\mathcal{A}, \mathcal{M})$ the Hochschild chain complex of
$\mathcal{A}$ with coefficients in $\mathcal{M}$ (a.k.a.~the cyclic
bar complex). Recall that this chain complex decomposes into a direct
sum
$CC(\mathcal{A}, \mathcal{M}) = \oplus_{d \geq 0} CC(\mathcal{A},
\mathcal{M}; d)$, with
$$CC(\mathcal{A}, \mathcal{M}; d) = \bigoplus_{\vec{X}}
\mathcal{M}(X_d,X_0) \otimes_{\k} \mathcal{A}(\vec{X}) ,$$ where the
direct sum runs over all tuples $\vec{X} = (X_0, \ldots, X_d)$ of
objects in $\mathcal{A}$. For $d=0$ we set $\mathcal{A}(\vec{X}) = \k$
so that
$CC(\mathcal{A}, \mathcal{M}; 0) := \bigoplus_{X \in \Ob(\mathcal{A})}
\mathcal{M}(X,X)$. We endow $CC(\mathcal{A}, \mathcal{M})$ with
the Hochschild differential $d_{CC}$ (see e.g.~\cite{Gan:thesis,
  Sher:form-hodge, Ri-Sm:WF-OC}). Its homology is denoted by
$HH(\mathcal{A}, \mathcal{M})$ and is called the Hochschild homology
of $\mathcal{A}$ with coefficients in $\mathcal{M}$.

In case $\mathcal{M} = \Delta_{\mathcal{A}}$ is the diagonal
$(\mathcal{A}, \mathcal{A})$-bimodule we write
$CC(\mathcal{A}, \mathcal{A})$ and $HH(\mathcal{A}, \mathcal{A})$ for
the Hochschild chain complex and homology respectively, with
coefficients in $\Delta_{\mathcal{A}}$. Sometimes we even abbreviate these to $CC(\mathcal{A})$ and $HH(\mathcal{A})$.

Assume now that $\mathcal{A}$ is a filtered $A_{\infty}$-category and
$\mathcal{M}$ is a filtered $(\mathcal{A}, \mathcal{A})$-bimodule.
For every $\alpha \in \mathbb{R}$, $d \geq 0$, and
$\vec{X} = (X_0, \ldots, X_d)$ denote by
$$( \mathcal{M}(X_d, X_0) \otimes \mathcal{A}(\vec{X}) )^{\alpha}
\subset \mathcal{M}(X_d,X_0) \otimes_{\k} \mathcal{A}(\vec{X})$$ the
$\k_0$-submodule generated by the images of all the submodules
$$\mathcal{M}^{\alpha'}(X_d, X_0) \otimes_{\k_0}
\mathcal{A}^{\alpha''_1}(X_0, X_1) \otimes_{\k_0} \cdots
\otimes_{\k_0}\mathcal{A}^{\alpha''_d}(X_{d-1}, X_d) , \text{ with }
\alpha' + \alpha''_1 + \cdots + \alpha''_d \leq \alpha,$$ under the
obvious map
$$\mathcal{M}(X_d, X_0) \otimes_{\k_0} \mathcal{A}(X_0, X_1) \otimes_{\k_0} \cdots
\otimes_{\k_0}\mathcal{A}(X_{d-1}, X_d) \longrightarrow
\mathcal{M}(X_d, X_0) \otimes_{\k} \mathcal{A}(\vec{X}).$$ Denote
$$CC^{\alpha}(\mathcal{A}, \mathcal{M}; d) := \bigoplus_{\vec{X}}
( \mathcal{M}(X_d,X_0) \otimes \mathcal{A}(\vec{X}) )^{\alpha}$$ which
is a $\k_0$-submodule of $CC(\mathcal{A}, \mathcal{M}; d)$. Finally,
denote by
$$CC^{\alpha}(\mathcal{A}, \mathcal{M}) := \bigoplus_{d \geq 0}
CC^{\alpha}(\mathcal{A}, \mathcal{M}; d).$$ This gives us an
increasing filtration of $CC(\mathcal{A}, \mathcal{M})$ by
$\k_0$-submodules
$CC^{\alpha}(\mathcal{A}, \mathcal{M}) \subset CC(\mathcal{A},
\mathcal{M})$, $\alpha \in \mathbb{R}$. Since $\mathcal{A}$,
$\mathcal{M}$ are filtered the Hochschild differential preserves this
filtration, i.e.
$$d_{CC}(CC^{\alpha}(\mathcal{A}, \mathcal{M})) \subset
CC^{\alpha}(\mathcal{A}, \mathcal{M})$$ so
$(CC(\mathcal{A}, \mathcal{M}), d_{CC})$ becomes a filtered chain
complex. We denote by $PHH(\mathcal{A}, \mathcal{M})$ its persistence
homology.

The Hochschild complex $CC(\mathcal{A},\mathcal{M})$ admits also a discrete filtration, traditionally called the \textit{length} filtration, $(F^NCC(\mathcal{A},\mathcal{M}))_{N\geq 0}$ defined for $N\geq 0$ as $$F^NCC(\mathcal{A},\mathcal{M}):=\bigoplus_{d\leq N}CC(\mathcal{A},\mathcal{M};d).$$ It is easy to see that $d_{CC}$ preserves this filtration too, and hence induces a differential on each $F^NCC(\mathcal{A},\mathcal{M})$. We denote by $F^HH(\mathcal{A},\mathcal{M})$ the resulting homology. The real filtration defined above is obviously compatible with the length filtrations. We thus get a sequence $F^NPHH(\mathcal{A},\mathcal{M})$, $N\geq 0$, of persistence modules.

Note that if $\mathcal{A}$ is filtered then the diagonal bimodule
$\Delta_{\mathcal{A}}$ is also filtered and we denote the
corresponding filtered chain complex and its persistence homology by
$CC(\mathcal{A}, \mathcal{A})$ and $PHH(\mathcal{A}, \mathcal{A})$
respectively.

We briefly describe the grading formalism for Hochschild homology used in our geometric application in \S\ref{sec:split-app}. Let $(X,\omega)$ be a monotone symplectic manifold of real dimension $2n$ and $\mathcal{A}$ be a Fukaya category build from Lagrangians in $(X,\omega)$. Floer complexes are defined over the Novikov field $\k=\La$ and are $\mathbb{R}/2\mathbb{Z}$ graded, with the Novikov variable having degree $0$. In particular, the following has to be understood $\mod 2$. The $d$th order $A_\infty$-map $\mu_d$ of $\mathcal{A}$ has cohomological degree $2-d$, hence homological degree $(d-2)+(d-1)n$. This leads to the following grading on the Hochschild complex $CC_*(\mathcal{A},\mathcal{A})$ of $\mathcal{A}$: given $\vec{x}:=x_1\otimes \cdots \otimes x_k\in \mathcal{A}(X_0,\ldots,X_k,X_0)$ we set $$\deg(\vec{x}):= \sum_{i=1}^k \deg(x_i)+k-1-n.$$ It is easy to see that $d_{CC}$ has degree $-1$.
\subsection{Shifted categories} \label{a:shifted} Let $\mathcal{A}$ be
a filtered $A_{\infty}$-category with a shift functor $\Sigma$, and
let $r \geq 0$. Define the normalized $r$-shift $S^r \mathcal{A}$ of
$\mathcal{A}$ to be the filtered $A_{\infty}$-category with the same
objects as $\mathcal{A}$ and morphism spaces $(S^r\mathcal{A})(L,L')$
between two objects $L$ and $L'$:
$$(S^r \mathcal{A})^{\alpha}(L,L'):=
\begin{cases} \mathcal{A}^{\alpha}(L,L'), \text{ if } L=\Sigma^lL'
  \text{ for some } l\in \mathbb R, \\ \mathcal{A}^{\alpha - r}
  (L,L'), \text{ otherwise.}
\end{cases}
$$
The $A_{\infty}$-operations $\mu^{\mathcal{A}}_d$ on $\mathcal{A}$
naturally induce maps $\mu^{S^r \mathcal{A}}_d$ on $S^r
\mathcal{A}$. The shift and translation functors $\Sigma$ and $T$
carry over to $S^r \mathcal{A}$ in the obvious way.

There is an obvious filtered $A_\infty$-functor
$\eta_r^{\mathcal{A}} : S^r \mathcal{A} \longrightarrow \mathcal{A}$
which is the identity on objects, is induced by the identity on
morphisms, and has trivial higher operations.



\subsection{Functors with Linear Deviation} \label{a:func-LD}

For readability, we repeat here the definition of
$A_{\infty}$-functors with linear deviation (LD functors, for short).

Let $\mathcal{A}$ be a filtered $A_{\infty}$-category.  Given a tuple
$\vec{X} = (X_0, \ldots, X_d)$ of objects from $\mathcal{A}$ we define
its reduced tuple $\vec{X}_R := (X_{i_0}, \ldots, X_{i_{d_R}})$ by
omitting from $\vec{X}$ subsequent (in the cyclic order) objects that
are equal up to a shift. The objects forming $\vec{X}_R$ are well
defined only up to shifts, but the length of $\vec{X}_R$ (or rather, its length-$1$)
$0 \leq d_R \leq d$ is well defined. We call it the reduced length of
$\vec{X}$ and denote it by $d_R$ or $d_R(\vec{X})$ whenever we want to
emphasize its dependence on $\vec{X}$. Note that in case every two
consecutive objects in $\vec{X}$ are different (up to shifts) then
$d_R(\vec{X}) = d$. At the other extremity, if all the objects in
$\vec{X}$ are equal up to shifts, then $d_R(\vec{X}) = 0$.

Let $\mathcal{A}$, $\mathcal{B}$ be two filtered
$A_{\infty}$-categories. Let $s\geq 0$. An $A_{\infty}$-functor
$\mathcal{F}: \mathcal{A} \longrightarrow \mathcal{B}$ is said to have
linear deviation rate $s$ if for every $d \geq 1$, every tuple
of objects $\vec{X} = (X_0, \ldots, X_d)$ from $\Ob(\mathcal{A})$ and
every $\alpha \in \mathbb{R}$ we have:
$$\mathcal{F}_d \bigl(\mathcal{A}^{\alpha}(X_0, \ldots, X_d)\bigr)
\subset \mathcal{B}^{\alpha + d_R(\vec{X}) s}(\mathcal{F}X_0,
\mathcal{F}X_d).$$ We will refer to such functors as LD-functors (LD
stands for Linear Deviation).

Let
$\mathcal{F}, \mathcal{G}: \mathcal{A} \longrightarrow \mathcal{B}$ be
two LD-functors, both with deviation rate $\leq s$.  A pre-natural
transformation $T: \mathcal{F} \longrightarrow \mathcal{G}$ is said to
have linear deviation rate $s$ and shift $\leq r$ if for every
$d \geq 0$ it's $d$-th order component $T_d$ shifts filtration levels
by $\leq r + ds$.

LD-functors $\mathcal{F}: \mathcal{A} \longrightarrow \mathcal{B}$
with a given deviation rate $s \geq 0$, and their pre-natural
transformations form a filtered $A_{\infty}$-category
$\text{fun}^{\text{LD};s}(\mathcal{A}, \mathcal{B})$. To emphasize the
given deviation rate we will sometimes denote the objects of these
category by $(\mathcal{F}, s)$. Note that for every $s' \leq s''$ we
have an obvious (faithful but not full) inclusion
$\text{fun}^{\text{LD};s'}(\mathcal{A}, \mathcal{B}) \subset
\text{fun}^{\text{LD};s''}(\mathcal{A}, \mathcal{B})$ of filtered
$A_{\infty}$-categories.

\subsubsection{Homotopy between LD-functors} \label{a:homotopy} We
also have the notion of homotopy between LD-functors, which is an
adaptation of the notion
from~\cite[Section~(1h)]{Se:book-fukaya-categ} to the filtered and LD
setting. Let
$\mathcal{F}, \mathcal{G}: \mathcal{A} \longrightarrow \mathcal{B}$ be
two LD-functors with deviation $\leq s$ and assume that both functors
{\em act in the same way on objects}.  Consider the pre-natural
transformation $D = \{D_d\}_{d \geq 0}$ between $\mathcal{F}$ and
$\mathcal{G}$ defined by $D_0 := 0$ and
$D_d := \mathcal{F}_d - \mathcal{G}_d$. We say that $\mathcal{F}$ and
$\mathcal{G}$ are $r$-homotopic, as LD-functors with deviation
$\leq s$ if there exists a pre-natural transformation
$T \in hom^r_{\text{fun}^{\text{LD};s}(\mathcal{A}, \mathcal{B})}\bigl(
(\mathcal{F},s), (\mathcal{G}, s) \bigr)$ with vanishing $0$-term,
$T^0=0$, and with deviation rate $\leq s$ and shift $\leq r$, such
that $D = \mu_1(T)$.

Recall that filtered functors are a special case of LD-functors, and
homotopies between them (as defined above; i.e.~with deviation $s=0$)
are interesting too. In this paper we will need homotopies both
between filtered functors as well as between LD-functors.

Next we briefly discuss the relation between homotopy of functors and
persistence Hochschild homology. Recall from \S\ref{app:PHH}, that if $\mathcal{A}$ is a filtered
$A_{\infty}$-category, then its Hochschild chain complex
$CC(\mathcal{A}, \mathcal{A})$ inherits a filtration from
$\mathcal{A}$ and consequently the homology of the latter becomes a
persistence module which we call the persistence Hochschild homology
$PHH(\mathcal{A}, \mathcal{A})$ of $\mathcal{A}$. Note that its
$\infty$-limit $PHH^{\infty}(\mathcal{A}, \mathcal{A})$ coincides with
the usual Hochschild homology $HH(\mathcal{A}, \mathcal{A})$.
\begin{prop} \label{p:homotopy-HH} Let
  $\mathcal{F}, \mathcal{G}: \mathcal{A} \longrightarrow \mathcal{B}$
  be two filtered functors. If $\mathcal{F}, \mathcal{G}$ are
  $0$-homotopic then the chain maps induced by $\mathcal{F}$,
  $\mathcal{G}$ on the filtered Hochschild chain complexes
  $\mathcal{F}_{\scriptscriptstyle CC},
  \mathcal{G}_{\scriptscriptstyle CC}:CC(\mathcal{A}, \mathcal{A})
  \longrightarrow CC(\mathcal{B}, \mathcal{B})$ are $0$-chain
  homotopic. In particular the induced maps on the persistence
  Hochschild homologies
  $\mathcal{F}_{\scriptscriptstyle PHH},
  \mathcal{G}_{\scriptscriptstyle PHH}: PHH(\mathcal{A}, \mathcal{A})
  \longrightarrow PHH(\mathcal{B}, \mathcal{B})$ coincide as maps of
  persistence modules.
\end{prop}

\subsubsection{From LD-functors to filtered functors} \label{a:ld-fil}

Let $\mathcal{F} : \mathcal{A} \longrightarrow \mathcal{B}$ be an LD
$A_{\infty}$-functor with deviation rate $\leq r$. Using \S\ref{a:shifted}, define a new
$A_{\infty}$-functor
$$\eta^r \mathcal{F}: S^r \mathcal{A} \longrightarrow \mathcal{B},
\quad \eta^r \mathcal{F} :=\mathcal{F} \circ \eta^{\mathcal{A}}_r,$$
which we call the normalized $r$-shift of $\mathcal{F}$. It is
straightforward to see that $\eta^r \mathcal{F}$ is a {\em filtered}
$A_{\infty}$-functor.



\subsection{Filtered bimodules associated with
  functors} \label{a:func-fil-bimod} Given a filtered
$A_{\infty}$-functor
$\mathcal{F}: \mathcal{A} \longrightarrow \mathcal{B}$, define the
$(\mathcal{B},\mathcal{A})$-bimodule
$$\underline{\mathcal{F}}:=
(\text{id}_\mathcal{B} \otimes \mathcal{F})^*\Delta_{\mathcal{B}},$$
where $\Delta_{\mathcal{B}}$ stands for the diagonal bimodule of
$\mathcal{B}$. In other words, given $Y \in \Ob(\mathcal{B})$,
$X \in \Ob(\mathcal{A})$, we have
$$\underline{\mathcal{F}}(Y,X):= \mathcal{B}(Y,\mathcal{F}(X)).$$
On bimodule-composable elements, the structure maps
$\mu^{\underline{\mathcal{F}}}_{l|1|r}$, for $l,r\geq 0$, of
$\underline{\mathcal{F}}$ are given by
\begin{equation*}
  \begin{aligned}
    & \mu^{\underline{\mathcal{F}}}_{l|1|r}(b_1,\ldots, b_l,y,a_1,\ldots, a_r)
    := \\
  & \sum_{j=1}^r \sum_{\substack{s_1 + \cdots + s_j = r \\ 1 \leq s_i}}
\mu^{\mathcal{B}}_{l+j+1}(b_1,\ldots, b_l, y,
\mathcal{F}_{s_1}(a_1,\ldots, a_{s_1}),\ldots,
    \mathcal{F}_{s_j}(a_{r-s_j},\ldots, a_r)).
  \end{aligned}
\end{equation*}
Note that since $\mathcal{F}$ is filtered, the bimodule
$\underline{\mathcal{F}}$ is filtered too.

The assignment $\mathcal{F} \longmapsto \underline{\mathcal{F}}$
extends to a filtered $A_{\infty}$-functor
$\mathcal{U}: F\text{fun}(\mathcal{A}, \mathcal{B}) \longrightarrow
F\text{bimod}_{\mathcal{B}, \mathcal{A}}$. Its 1st order term
$\mathcal{U}_1$ has the following description. Let
$\mathcal{F}, \mathcal{G}; \mathcal{A} \longrightarrow \mathcal{B}$ be
two filtered $A_{\infty}$-functors and
$T: \mathcal{F} \longrightarrow \mathcal{G}$ a pre-natural
transformation between them, of some given shift $\alpha$
(i.e.~$T \in \hom^{\alpha}_{F\text{fun}(\mathcal{A},
  \mathcal{B})}(\mathcal{F}, \mathcal{G})$). Define the
$(\mathcal{B}, \mathcal{A})$-bimodule pre-homomorphism
\begin{equation*}
  \begin{aligned}
    & \mathcal{U}_1(T)\colon \underline{\mathcal{F}} \longrightarrow \underline{\mathcal{G}},
    \\
    & \mathcal{U}_1(T)_{l|1|r}(b_1,\ldots, b_l, y,a_1,\ldots, a_r):= \\
    & \sum_{j=1}^r\sum_{i=1}^j\sum_{\substack
      {s_1 + \cdots + s_j = r  \\ 1 \leq s_k, \, 
    \forall k \neq i}}
    \mu^{\mathcal{B}}_{l+1+j}\bigl(b_1,\ldots, b_l,y, (\mathcal{F}_{s_1},\ldots,
    \mathcal{F}_{s_{i-1}},T_{s_i},\mathcal{G}_{s_{i+1}}, \ldots,
    \mathcal{G}_{s_j})(a_1,\ldots, a_r)\bigr).
  \end{aligned}
\end{equation*}
It is straightforward to verify that $\mathcal{U}_1(T)$ is bimodule
pre-homomorphism of the same shift $\alpha$. Given composable natural
transformations $T_1,\ldots, T_d$ between $A_\infty$-functors from
$\mathcal{A}$ to $\mathcal{B}$, the pre-homomorphism of
$(\mathcal{B},\mathcal{A})$-bimodules $\mathcal{U}_d(T_1,\ldots, T_d)$
is defined by a similar formula.

\begin{lem}[Theorem~8.2.1.2 in \cite{Lfv:thesis}]
  The component $\mathcal{U}_1$ of the functor $\mathcal{U}$ induces a
  quasi-isomorphism.
\end{lem}


\subsection{Pull back and push forward of filtered
  $A_{\infty}$-modules} \label{ap:pshf}

Let $\mathcal{A}$, $\mathcal{B}$ be filtered $A_{\infty}$-categories
and $\mathcal{F} : \mathcal{A} \longrightarrow \mathcal{B}$ a filtered
$A_{\infty}$-functor. Let $\mathcal{M}$ be a filtered
$\mathcal{B}$-module. Then the pull back $\mathcal{F}^*\mathcal{M}$
can be endowed with the structure of a filtered $\mathcal{A}$-module
in a straightforward way. However, if instead of assuming that
$\mathcal{F}$ is filtered we only assume it to be an LD-functor then
the pull back $\mathcal{F}^*\mathcal{M}$ is in general not filtered,
but rather an $LD$-module (a concept which will not be used in this
paper).

At the same time, there is a way to define push-forward of filtered
modules by LD-functors and still get filtered modules. We present this
construction next.

Let $\mathcal{A}, \mathcal{B}$ be filtered $A_{\infty}$-categories and
$\mathcal{F}: \mathcal{A} \longrightarrow \mathcal{B}$ an LD-functor
with deviation rate $\leq r$. Let $\mathcal{M}$ be a filtered
$\mathcal{A}$-module. The push-forward $(\mathcal{F},r)_* \mathcal{M}$
of $\mathcal{M}$ by $\mathcal{F}$ is defined to be the filtered
$\mathcal{B}$-module:
\begin{equation} \label{eq:psh-m} (\mathcal{F},r)_* \mathcal{M} :=
  \underline{\eta^r \mathcal{F}} \otimes_{S^r \mathcal{A}}
  (\eta_r^{\mathcal{A}})^* \mathcal{M}.
\end{equation}
Note that the value of the parameter $r$ affects the filtration
structure on the push-forward module, and therefore we include it in
the notation. In other words, if $r < r'$ then
$(\mathcal{F},r')_* \mathcal{M}$ has a different filtration structure
than $(\mathcal{F},r)_* \mathcal{M}$. Of course, if one forgets the
filtrations, then the push-forward of $\mathcal{M}$ depends only on
$\mathcal{F}$, and gives the same result for all $r$'s.

Given a pre-homomorphism
$f : \mathcal{M}_0 \longrightarrow \mathcal{M}_1$ beetween filtered
$\mathcal{A}$-modules, define its pushforward
$(\mathcal{F},r)_* f : (\mathcal{F},r)_* \mathcal{M}_0 \longrightarrow
(\mathcal{F},r)_* \mathcal{M}_1$ as follows. The higher order terms
$((\mathcal{F}, r)_*f)_{l|1}$, $l \geq 1$, are defined to be $0$, and
the linear order term is:
$$((\mathcal{F}, r)_*f)_{0|1}
(b\otimes a_1\otimes \cdots \otimes a_d\otimes m)= \sum_{i=1}^d
b\otimes a_1\otimes \cdots \otimes a_i\otimes
f_{d-i|1}(a_{i+1},\ldots, a_d,m).$$

The construction above extends to a filtered $A_{\infty}$-functor
$$PF: \text{fun}^{\text{LD};r}(\mathcal{A}, \mathcal{B})
\longrightarrow F\text{fun}(F\md_{\mathcal{A}}, F\md_{\mathcal{B}}),$$
whose action on the object $(\mathcal{F}, r)$ is the filtered
functor
$(\mathcal{F},r)_* :F\md_{\mathcal{A}} \longrightarrow
F\md_{\mathcal{B}}$. To define the action of $PF$ on morphisms, let
$\mathcal{F}, \mathcal{G} \in \Ob
(\text{fun}^{\text{LD};r}(\mathcal{A}, \mathcal{B}))$, and let
$T\colon \mathcal{F}\to \mathcal{G}$ be a pre-natural transformation
(in the category $\text{fun}^{\text{LD};r}(\mathcal{A}, \mathcal{B})$
of LD-functors with deviation rate $\leq r$) of shift $\rho \geq
0$. The image $PF_1(T)$ of $T$ under the push-forward functor $PF$ is
defined for an object $\mathcal{M} \in \Ob(F\md_{\mathcal{A}})$ as the
pre-module homomorphism
$PF_1(T) \in \hom_{F\md_{\mathcal{B}}}(\mathcal{F}_*\mathcal{M},
\mathcal{G}_*\mathcal{M})$ given by:
\begin{align*}
  PF(T)_{l|1}(b_1,\ldots, b_l, y\otimes \vec{a}\otimes m) & :=
  \sum_{j=0}^d \mathcal{U}_1(T)_{l|1|j}(b_1,\ldots, b_l,y,a_1,\ldots,
  b_j)\otimes a_{j+1}\otimes \cdots \otimes a_d\otimes m \\ =
  \sum_{j=0}^d \sum_{l=1}^j\sum_{i=1}^l\sum_{\substack{s_1 + \cdots +
      s_l = j \\ 1 \leq s_k, \, \forall k \neq i}}
  \mu^{\mathcal{B}}_{l+1+j}\bigl( & b_1,\ldots, b_l,y,(\mathcal{F}_{s_1},\ldots,
  \mathcal{F}_{s_{i-1}},T_{s_i},\mathcal{G}_{s_{i+1}}, \ldots,
  \mathcal{G}_{s_l})(a_1,\ldots, a_j)\bigr)  \\
  & \otimes a_{j+1}\otimes
  \cdots \otimes a_d\otimes m.
\end{align*}
The image of a composable tuple $T_1,\ldots, T_d$ of
$A_\infty$-pre-natural transformations under $PF_d$ is defined in a
very similar manner. It is straightforward to verify that $PF$ is a
filtered (strictly unital) $A_{\infty}$-functor.

The next proposition shows that, up to shifts, the push forward of the
Yoneda module of an objects is the same as the Yoneda module of the
image of this object by the given functor. We will denote here by
$\mathcal{Y}$ the filtered Yoneda embedding, both for $\mathcal{A}$
and for $\mathcal{B}$.
\begin{prop} \label{p:pushf-yon} Let
  $\mathcal{F}: \mathcal{A} \longrightarrow \mathcal{B}$ be an
  LD-functor with deviation rate $\leq r$. For every
  $X \in \Ob(\mathcal{A})$ the push forward
  $(\mathcal{F},r)_* \mathcal{Y}(X)$ of the Yoneda module
  $\mathcal{Y}(X)$ is $0$-quasi-isomorphic to
  $\Sigma^r \mathcal{Y}(\mathcal{F} X)$.
\end{prop}

To simplify the notation, when the deviation rate of a functor is
clear from the context, we will sometimes omit the $r$ from the pair
$(\mathcal{F}, r)$ and simply write $\mathcal{F}_* \mathcal{M}$.

\subsection{Filtered twisted complexes}\label{ap:ftc} In this section, we introduce filtered twisted complexes. The constructions developed here are upgrades of \cite[Chapter 3]{Se:book-fukaya-categ} from the unfiltered to the filtered setting, and of \cite[Section 2.5.1]{BCZ:tpc} from the $dg$ to the $A_\infty$ case.
\subsubsection{The filtered $A_\infty$-categories of filtered twisted complexes} Let $\mathcal{A}$ be a filtered $A_\infty$-category, without a choice of shift functor.
We will write $\k_0^r[d]$ for interval persistence module over $\k_0$ with a unique generator in degree $d\in \mathbb Z$ placed at filtration level $r\in \mathbb R$ and given an object $L$ of $\mathcal{A}$ write $\Sigma^rL[d]$ for the formal tensor product $\k_0^r[d]\otimes L$. 
\begin{dfn}
	We define the (standard) shift completion of $\mathcal{A}$ as the category $\mathcal{A}^\Sigma$ with objects $$\text{Ob}(\mathcal{A}^\Sigma) := \left\{ \Sigma^rL[d]\ : \ L\in \text{Ob}(\mathcal{A}), \ d\in \mathbb{Z}, \ r\in \mathbb R \right\},$$ morphism spaces $$\mathcal{A}^\Sigma(\Sigma^rL_0[d_0],\Sigma^sL_1[d_1]) := \Sigma^{-(r-s)}\mathcal{A}(L_0,L_1)[d_1-d_0]$$ that is, $\mathcal{A}(\Sigma^rL_0[d_0],\Sigma^sL_1[d_1]) $ is the chain complex $\mathcal{A}(L_0,L_1)$ shifted up in degree by $d_1-d_0$ and in filtration by $r-s$, and with the same $A_\infty$-operations of $\mathcal{A}$.
\end{dfn}
\begin{rem}
	\begin{enumerate}
		\item It is easy to see that $\mathcal{A}^\Sigma$ is a filtered $A_\infty$-category. 
		\item On the $A_\infty$-category $\mathcal{A}^\Sigma$ there is an immediate choice of shift $A_\infty$-functor (in the sense of \cite[Section 3.2.1]{BCZ:tpc}), which on objects is $\Sigma^r(L) := \Sigma^rL$. %Moreover, this shift functor is compatible in the sense of Definition \ref{compatibilityofshift}.
		\item In some texts the prefix $\Sigma$ denotes the additive enlargement of an $A_\infty$-category (e.g. \cite[Section (3k)]{Se:book-fukaya-categ}), which is not the case here.
	\end{enumerate}
\end{rem}
\begin{dfn}
	We define the filtered additive enlargement $\mathcal{A}^\oplus$ of $\mathcal{A}$ as the additive enlargement of $\mathcal{A}^\Sigma$ in the sense of \cite[Section (3k)]{Se:book-fukaya-categ}, that is, as the $A_\infty$-category with objects given by formal sums $$\overline{L} = \bigoplus_{i=1}^n\Sigma^{r_i}L_i[d_i],$$ where $n\in \mathbb N$ and $\Sigma^{r_i}L_i[d_i]$ is an object of $\mathcal{A}^\Sigma$ for any $i$; and morphisms spaces between objects $\overline{L}=\bigoplus_{i=1}^n\Sigma^{r_i}L_i[d_i]$ and $\overline{L}'=\bigoplus_{i=1}^m\Sigma^{r_i'}L'_i[d'_i]$ given by $n\times m$ matrices $$f = (f_{ij}), \text{  where } f_{ij}\in \mathcal{A}^\Sigma\left(\Sigma^{r_i}L_i[d_i],\Sigma^{r_j'}L_j'[d_j']\right)$$ with usual matrix grading. Given an object $\overline{L}=\bigoplus_{i=1}^n\Sigma^{r_i}L_i[d_i]$ we will write $|\overline{L}|=n$ and call it the length of $\overline{L}$. The $A_\infty$-maps $\mu_d^\oplus$, $d\geq 1$, are defined by extending $\mu_d$ matrixwise, that is, given $\overline{L}_0,\ldots, \overline{L}_d$ objects of $\mathcal{A}^\oplus$ and matrices $f^1,\ldots,f^d$ with $f^i\in \mathcal{A}^\oplus(\overline{L}_{i-1},\overline{L}_i)$ we set for $i=1,\ldots, |\overline{L}_0|$ and $j=1,\ldots, |\overline{L}_d|$: $$\mu_d^\oplus\left(f^1,\ldots, f^d\right)_{ij}:= \sum_{i_1=1}^{|\overline{L}_1|}\cdots \sum_{i_{d-1}=1}^{|\overline{L}_{d-1}|}\mu_d\left(f^1_{ii_1},\ldots, f^d_{i_{d-1}j}\right).$$
	Moreover, we filter morphism spaces in $\mathcal{A}^\oplus$ by setting the filtration level of a matrix to be the maximum of the filtration level of its entries, that is: $$(\mathcal{A}^\oplus)^{\leq \alpha}(\overline{L},\overline{L}') := \left\{f= (f_{ij})\in \mathcal{A}^\oplus(\overline{L},\overline{L}'):\ f_{ij}\in (\mathcal{A}^\Sigma)^{\leq \alpha}\left(\Sigma^{r_i}L_i[d_i], \Sigma^{r_j'}L_j'[d_j']\right )\text{  for any }i,j\right\}.$$
\end{dfn} 

\begin{rem}
	\begin{enumerate}
	%	\item In the notation of \cite{Se:book-fukaya-categ} we would write $\mathcal{A}^\oplus$ as $\Sigma \mathcal{A}^\Sigma$.
		\item We will not write all the sums in the definition of $\mu_d^\oplus$, but rather use the shorthand notation $\sum_{i_1,\ldots, i_{d-1}}$.
		\item It is straightforward to see that $\mathcal{A}^\oplus$ is a filtered $A_\infty$-category.
	\end{enumerate}
\end{rem}

\begin{dfn}\label{defFTW}
	A filtered (one-sided) twisted complex of $\mathcal{A}^\Sigma$ is a pair $$(\overline{L}, q=q_{\overline{L}}):= \left(\bigoplus_{i=1}^n \Sigma^{r_i}L_i[d_i], (q_{ij})\right)$$ for some $n\in \mathbb N$ such that:
	\begin{enumerate}
		\item $\overline{L}$ is an object of $\mathcal{A}^\oplus$, that is, for any $i=1,\ldots, n$, $ \Sigma^{r_i}L_i[d_i]$ is an object of $\mathcal{A}^\Sigma$,
		\item $q$ is a morphism in $\mathcal{A}^\oplus (\overline{L}, \overline{L})$ lying at degree $1$ and vanishing filtration level.
		\item the matrix $q$ is strictly upper triangular, that is $q_{ij}=0$ for any $i\geq j$,
		\item\label{cond4} the matrix $q$ satisfies $$\sum_{d\geq 1}\mu_d^\oplus (q,\ldots, q)=0,$$ 
		%which we cafor any $i,j=1,\ldots, n$ we have $$\sum_{d\geq 1}\sum_{i_1,\ldots ,i_{d-1}=1}^n\mu_d(q_{ii_1}, q_{i_1i_2},\ldots , q_{i_{d-1}j}) = 0,$$
		to which we will refer as the Maurer-Cartan identity of $q$.
	\end{enumerate}
	The matrix $q=q_{\overline{L}}$ is called the differential of $\overline{L}$ and will be often dropped from the notation.
\end{dfn}
\begin{rem} As $q$ is strictly upper triangular, any entry $\mu_d^\oplus(q,\ldots, q)_{i,j}$, $i,j=1,\ldots, |\overline{L}|$, can be written as $$\sum_{i<i_1<\cdots < i_{d-1}<j} \mu_d(q_{ii_1}, q_{i_1i_2},\ldots , q_{i_{d-1}j}).$$ Moreover, the sum in the Maurer-Cartan identity of $q$ is a finite sum: indeed, for any $d\geq |\overline{L}|$ we have $\mu_d^\oplus(q,\ldots, q)=0$ since the longest possible term is $$\mu_{|L|-1}(q_{12},q_{23},\ldots, q_{|\overline{L}|-1|\overline{L}|})$$ in the matrix $\mu^\oplus_{|L|-1}(q,\ldots, q)$, which has only one non-zero entry in position $(1,|\overline{L}|)$.
\end{rem}
\noindent Given a twisted complex over a filtered $A_\infty$-category in the sense of \cite[Section (3l) and Remark 3.26]{Se:book-fukaya-categ}, it might be the case that the differential is not filtration-preserving (i.e. Condition \ref{cond4} above does not hold). The analogue of Lemma 2.96 in \cite{BCZ:tpc} holds: there are shifts $r_i$ turning this twisted complex in a filtered one.
\begin{dfn}
	We define the filtered pre-triangulated completion $FTw(\mathcal{A})$ of $\mathcal{A}$ as follows:
	\begin{enumerate}
		\item The objects of $FTw(\mathcal{A})$ are filtered one-sided twisted complexes over $\mathcal{A}$,
		\item Given two filtered twisted complexes $\overline{L}$ and $\overline{L}'$, the morphism space $FTw(\mathcal{A})(\overline{L},\overline{L}')$ is defined as $\mathcal{A}^\oplus(\overline{L},\overline{L}')$, with the same filtration, %are matrices $f=(f_{ij})$ such that $$f_{ij} \in \mathcal{A}^\Sigma\left(\Sigma^{r_i}L_i[d_i], \Sigma^{r_j'}L_j'[d_j']\right),$$ 
		\item The $A_\infty$-operations are deformed by the differentiasl in the following way: given any $d\geq 1$ and twisted complexes $$\left(\overline{L_i} = \bigoplus_{k=1}^{n_i}\Sigma^{r_{i,k}}L_{i,k}[d_{i,k}],q_i\right)$$ for any $i=0,\ldots, d$ and morphisms $f^i\in FTw(\mathcal{A})(\overline{L_{i-1}},\overline{L_i})$ for any $i=1,\ldots, d$ we define $\mu_d^{Tw}(f^1,\ldots, f^d)$ as the $|\overline{L}_0|\times |\overline{L}_d|$ matrix  $$\mu_d^{Tw}(f^1,\ldots, f^d) := \sum_{k_0,\ldots, k_d\geq 0}\mu_{d+k_0+\cdots + k_d}^{\oplus}\left(q_0^{\otimes k_0}, f^1,q_1^{\otimes k_1},\ldots, q_{d-1}^{\otimes k_{d-1}},f^d, q_d^{\otimes k_d}\right)$$
		%for $i=1,\ldots, n_0$ and $j=1,\ldots, n_d$ the element $$\mu_d^{FTw(\mathcal{A})}(f^1,\ldots, f^d)_{ij} := \sum_{k_0,\ldots, k_d\geq 0} \mu_{d+k_0+\cdots + k_d}^{\oplus}\left(q_0^{\otimes k_0}, f^1,q_1^{\otimes k_1},\ldots, q_{d-1}^{\otimes k_{d-1}},f^d, q_d^{\otimes k_d}\right)$$ where $\mu_d^{\oplus}$ is the matrixwise extension of $\mu_d^\mathcal{A}$ to direct sums as in \cite[Remark 3.26]{Se:book-fukaya-categ}. {\color{orange} is this understandable?}
		%\item Given two filtered twisted complexes $\overline{L}$ and $\overline{L}'$ as above, the filtration on $FTw(\mathcal{A})(\overline{L},\overline{L}')$ is defined by taking the maximum in filtration of the entries of a morphism, i.e. $$FTw(\mathcal{A})^{\leq \alpha}(\overline{L},\overline{L}') := \left\{f= (f_{ij})\in FTw(\mathcal{A})(\overline{L},\overline{L}'):\ f_{ij}\in (\mathcal{A}^\Sigma)^{\leq \alpha}\left(\Sigma^{r_i}L_i[d_i], \Sigma^{r_j'}L_j'[d_j']\right )\text{  for any }i,j\right\}.$$
		%{\color{purple} here we are at the chain level, so we are dealing with inclusions indeed, and the above def makes sense (?)}
	\end{enumerate}
	Moreover, given $N\geq 1$ we define $FTw^N\mathcal{A}$ to be the full $A_\infty$-subcategory of $FTw\mathcal{A}$ with filtered twisted complexes $\overline{L}$ of length $|\overline{L}|\leq N$ as objects.
\end{dfn}
\begin{rem}\label{finitenesssum} 
a. The fact that $FTw(\mathcal{A})$ and $FTw^N(\mathcal{A})$ are a $A_\infty$-category is a consequence of differentials satisfying the Maurer-Cartan identity above.

b. There is an obvious filtered full and faithful $A_\infty$-functors $\mathcal{I}^N_\mathcal{A}\colon \mathcal{A}\to FTw^N(\mathcal{A})$ for all $N\geq 1$, as well as $\mathcal{I}_\mathcal{A}\colon \mathcal{A}\to FTw(\mathcal{A})$.

c. The shift functor $\Sigma$ on $\mathcal{A}^\Sigma$ induces a shift functor, still denoted by $\Sigma$, on $FTw(\mathcal{A})$. It is defined on an object $\overline{L}= \bigoplus_{i=1}^n \Sigma^{r_i}L_i[d_i]$ by $$\Sigma^r\overline{L} = \bigoplus_{i=1}^n \Sigma^{r_i+r}L_i[d_i]$$ and viewing each entry $q_{ij}$ of the differential $q$ of $\overline{L}$ as a morphism $$q_{ij}\in FTw(\mathcal{A})(\Sigma^{r_i+r}L_i[d_i], \Sigma^{r_j+r}L_j[d_j])\cong FTw(\mathcal{A})(\Sigma^{r_i}L_i[d_i], \Sigma^{r_j}L_j[d_j]).$$

d. The sum in the definition of $\mu_d^{Tw}$ is a finite sum, for basically the same reason that the Maurer-Cartan identity consists of a finite sum. For instance, in $\mu_1^{Tw}f$ for $f\in FTw(\mathcal{A})(\overline{L},\overline{L}')$ we sum elements of the form $\mu_{1+k_0+k_1}(q^{\otimes k_0},f,q'^{\otimes k_1})$, whose $(i,j)$-entry is $$\sum_{i_1,\ldots, i_{k_0}=1}^{|\overline{L}|}\sum_{i_{k_0+1},\ldots, i_{k_0+k_1}=1}^{|\overline{L}'|}\mu_{1+k_0+k_1}\left(q_{ii_1},\ldots, q_{i_{k_0-1}i_{k_0}},f_{i_{k_0}i_{k_0+1}},q'_{i_{k_0+1}i_{k_0+2}},\ldots, q'_{i_{k_0+k_1}j}\right)$$ so that, by lower triangularity of the differentials, $\mu_{1+k_0+k_1}(q^{\otimes k_0},f,q'^{\otimes k_1})=0$ whenever $k_0\geq |\overline{L}|$ or $k_1\geq |\overline{L}'|$. In particular, the longest possible summand is $$\mu_{|\overline{L}|+|\overline{L}'|-1}\left(q_{12},\ldots, q_{|\overline{L}|-1|\overline{L}|},f_{|\overline{L}|1},q'_{12},\ldots, q'_{|\overline{L}'|-1|\overline{L}'|}\right)$$ in the $(1,|\overline{L}'|)$-entry. It is easy to see that the same holds in $\mu_d^{Tw}(f^1,\ldots, f^d)$ whenever $k_i\geq |\overline{L_i}|$ for some $i$.

\end{rem}

\begin{dfn}
	Let $\overline{L}$ and $\overline{L}'$ be two filtered twisted complexes in $FTw(\mathcal{A})$ as above, $f\in FTw(\mathcal{A})(\overline{L}, \overline{L}')$ be a degree zero morphism such that $\mu_1^{FTw(\mathcal{A})}f=0$, and $\lambda\in \mathbb R$ such that $\lambda\geq \mathbb{A}(f)$. We define the $\lambda$-filtered mapping cone of $f$ as $$\text{Cone}^{\lambda}(f) := \left(\overline{L}'\oplus \Sigma^{\lambda}\overline{L}[1], \ q_{co}\right) \text{   where  }q_{co} := \begin{pmatrix}
		q' & f \\
		0 & \Sigma^\lambda q[1]
	\end{pmatrix}.$$
\end{dfn}
\begin{lem}
	$\text{Cone}^{\lambda}(f)$ is a filtered twisted complex.
\end{lem}
\begin{rem}
	Note that the subcategories $FTw^N(\mathcal{A})$ are not pre-triangulated, as the cone construction doesn't preserve the length filtration.
\end{rem}
\begin{prop}
	Let $\mathcal{A}$ be a filtered $A_\infty$-category. Then $H^0(FTw(\mathcal{A}))$ is a TPC.
\end{prop}


\begin{rem}
	Filtered twisted complexes might be described by allowing tensoring with graded filtered vector spaces (not persistence modules!) which are not one dimensional. Given a finitely and freely generated filtered graded vector $V := \bigoplus_{i\in I}\Lambda\cdot \gamma_i$ over $\Lambda$ and an object $L$ of $\mathcal{A}$, we can define the tensor $V\otimes L$ simply as the direct sum (or in other words, $0$-filtered mapping cone of the zero maps) of the elements $\Sigma^{|\gamma_i|}L[\deg(\gamma_i)] = \Lambda\cdot \gamma_i\otimes L$, where $\deg(\gamma_i)$ denotes the degree of $\gamma_i$ in $V$. We will often denote $\Sigma^{|\gamma_i|}L[\deg(\gamma_i)]$ simply as $L\cdot \gamma_i$. If $V$ is a filtered chain complex, i.e. it carries a differential $d_V$, then the tensor product inherits a deformed differential as follows: write $d_V\gamma_i = \sum_j\alpha^i_j\gamma_j$ for any $i\in I$, then the differential on $V\otimes L$ is the matrix with $(i,j)$ entry equal to $\alpha^i_je_L$, where $e_L\in \mathcal{A}(L,L)$ is the strict unit. This definition extends to the whole category $FTw(\mathcal{A})$ (i.e. it's not only well-defined for images of objects of $\mathcal{A}$) in an obvious way.  Another possibility (cfr. \cite[Section (3l)]{Se:book-fukaya-categ}) is to define the category $FTw$ as having as objects tensors of the form $\bigoplus V_i\otimes L_i$, where $V_i$ is a finite dimensional filtered vector space and $L_i$ is an object of $\mathcal{A}$, together with an upper triangular differential $q$. In this case morphisms between $\bigoplus_i V_i\otimes L_i$ and $\bigoplus_j W_j\otimes L'_j$ are elements of $$\bigoplus_{i,j}\text{Lin}(V_i,W_j)\otimes \mathcal{A}(L_i,L_j')$$ and the $A_\infty$-maps are modified by considering compositions of linear maps (see \cite[Equation (3.17)]{Se:book-fukaya-categ}). It is striaghtforward to see that the two descriptions are equivalent.
\end{rem}

%=======================
%=======================
%=======================

%\subsection{Retract approximability}
\subsubsection{The filtered extended Yoneda embedding}\label{extendedYoneda}
We extend the Yoneda embedding $\mathcal{Y}_\mathcal{A}\colon \mathcal{A}\to F\md_{\mathcal{A}}$ introduced above to a filtered $A_\infty$-functor $$\widetilde{\mathcal{Y}_\mathcal{A}}\colon FTw\mathcal{A}\to F
\md_{\mathcal{A}}$$ via $$\widetilde{\mathcal{Y}_\mathcal{A}}:= \mathcal{I}_\mathcal{A}^*\circ \mathcal{Y}_{FTw\mathcal{A}}$$ where $\mathcal{I}_\mathcal{A}^*$ is the pullback of the $A_\infty$-functor $\mathcal{I}_\mathcal{A}\colon \mathcal{A}\to FTw(\mathcal{A})$ introduced in Remark \ref{finitenesssum}b., $\mathcal{Y}_{FTw\mathcal{A}}$ is the filtered Yoneda embedding for the filtered $A_\infty$-category $FTw\mathcal{A}$ and $\circ$ is the usual composition of $A_\infty$-functors. As pullbacks of filtered $A_\infty$-functors preserve filtered $A_\infty$-modules and are themselves filtered, the functor $\widetilde{\mathcal{Y}_A}$ is indeed a filtered $A_\infty$-functor. Note that this is just an immediate extension to the filtered world of the definition of the extended Yoneda embedding in \cite[Section (3s)]{Se:book-fukaya-categ}.
\begin{prop}
	The induced functor $H(\widetilde{\mathcal{Y}_\mathcal{A}})\colon PD(\mathcal{A})\to H(F	\md_{\mathcal{A}})$ gives a TPC equivalence between $PD(\mathcal{A})$ and the image of $H(\widetilde{\mathcal{Y}_\mathcal{A}})$.
\end{prop}
\begin{proof}
	The fact that $H(\widetilde{\mathcal{Y}_\mathcal{A}})$ gives an equivalence of persistence categories is clear. We need to prove that it is a TPC functor, that is, it is compatible with the shift functors and that it is triangulated at the $0$-level. The fact that it is compatible with shift functors is a direct consequence of the fact that the shift functor on $FTw\mathcal{A}$ is compatible with modules by construction. The fact that the $0$-level is triangulated follows from the Yoneda embedding being an $A_\infty$-functor at the chain level.
\end{proof}
%===============================
%===============================
%==============================
\subsubsection{Filtered twistings}\label{filtwistings}
%Let $Y$ be an object of our filtered $A_\infty$-category $\mathcal{A}$ and $\overline{L}=(L,q_L)$ be a filtered twisted complex over $\mathcal{A}$. We can consider the filtered twisted complex $$Y\otimes FTw(Y,\overline{L})$$ with differential 
Let $X$ and $Y$ be objects of $\mathcal{A}$, and consider the filtered twisted complex $Y\otimes \mathcal{A}(X,Y)$ with differential $$e_Y\otimes \mu_1\in FTw(Y\otimes \mathcal{A}(X,Y),Y\otimes \mathcal{A}(X,Y)) = \mathcal{A}(Y,Y)\otimes \text{Lin}(\mathcal{A}(X,Y), \mathcal{A}(X,Y))$$ where $e_Y\in \mathcal{A}(Y,Y)$ is the strict unit of $Y$ and $\text{Lin}(-,-)$ denotes the space of linear maps between two vector spaces. In other words, given a basis $\vec{a}=(a_1,\ldots, a_n)$ of $\mathcal{A}(Y,X)$, the twisted complex above is $\bigoplus_{i=1}^n \Sigma^{|a_i|}Y[\text{deg}(a_i)]$ with differential given by the matrix $q_{ij} = \alpha_{ij}e_Y$, where $e_Y\in \mathcal{A}(Y,Y)$ is the strict unit for $Y$ and the $\alpha_{ij}\in \k$ are the coeffients in $\mu_1(a_i)=\sum_{j=1}^n\alpha_{ij}a_j$. We define the filtered morphism of filtered twisted complexes $$\xi\colon Y\otimes \mathcal{A}(Y,X)\to X$$ to be the map corresponding to the identity under the canonical filtered isomorphism of chain complexes $$FTw(Y\otimes \mathcal{A}(Y,X), X)\cong \text{Lin}(\mathcal{A}(Y,X),\mathcal{A}(Y,X)).$$ 
\begin{dfn}
	We define the twisting $T_YX$ of $X$ by $Y$ as the twisted complex given by the cone of $\xi$, i.e. $T_YX = \text{Cone}(\xi)$
\end{dfn}
We can write $$T_YX= X\oplus \bigoplus_{i=1}^n \Sigma^{|a_i|}Y[\text{deg}(a_i)+1]$$ with differential given by the matrix $\begin{pmatrix}
	0& \vec{a} \\ \vec{0}^T & (q_{ij})
\end{pmatrix}$.
Note that the strict unit $e_X\in \mathcal{A}(X,X)$ induces a morphism of twisted complexes $i\colon X\to T_YX$ as a vector $(e_X,0,\ldots,0)$.

We are interested in describing the image of $T_Y X$ under the extended filtered Yoneda embedding $\widetilde{\mathcal{Y}_\mathcal{A}}$ introduced in \S\ref{extendedYoneda}.
\begin{lem}
	We have that $$\tilde{\mathcal{Y}}_{T_YX} \cong \text{Cone}(\phi)$$
	where $\phi\colon \mathcal{Y}_Y\otimes \mathcal{A}(Y,X)\to \mathcal{Y}_X$ is the filtered  morphism of modules given by $$\phi_{l|1}(x_1,\ldots, x_l,a\otimes b):= \mu_{s+2}(x_1,\ldots, a,b)$$ and $\mathcal{Y}_Y\otimes \mathcal{A}(Y,X)$ has a module structure as defined in \cite[Section (3c)]{Se:book-fukaya-categ}.
\end{lem}
\begin{proof}
	It is easy to see, that as an element of $FTw(Y\otimes \mathcal{A}(Y,X), X) = \mathcal{A}(Y,X)\otimes \text{Lin}(A(Y,X),\Lambda)$, the morphism $\xi$ corresponds to the element $$\sum_{i=1}^na_i\otimes \psi_{a_i}$$ where $a_i\in \mathcal{A}(Y,X)$ is a basis element as above, and $\psi_{a_i}\colon A(Y,X)\to \Lambda$ is the map sending $a_i$ to $1$ and all the other basis elements to $0$. Under the Yoneda embedding, i.e. as a morphism of $A_\infty$-modules $\mathcal{Y}_Y\otimes \mathcal{A}(Y,X)\to \mathcal{Y}_X$, $\xi$ becomes the map \begin{align*}
		\phi_{l|1}(x_1,\ldots, x_l,b\otimes c) & = \sum_{i=1}^n \psi_{a_i}(c)\lambda(a_i)_{l|1}(x_1,\ldots, x_l,b)\\ & = \sum_{i=1}^n \psi_{a_i}(c)\mu_{l+2}(x_1,\ldots, x_l,b,a_i) = \mu_{l+2}(x_1,\ldots, x_l,b,c).
	\end{align*}
	This ends the proof.
\end{proof}
We can easily extend this construction to the case where $X$ is itself a filtered twisted complex over $\mathcal{A}$. In this case, under the extended Yoneda embedding, $X$ corresponds to a filtered $A_\infty$-module $\mathcal{M}$ and $T_YX$ corresponds to the $0$-cone of the full contraction map $$\phi\colon \mathcal{Y}_Y\otimes \mathcal{M}(Y)\to \mathcal{M}$$ (cfr. \cite[Section (5a)]{Se:book-fukaya-categ}). We will write the above cone as $$\Tau_Y\mathcal{M}:= \text{Cone}(\phi).$$ The analogous of the lemma above holds in this case too.


\begin{lem}
	Let $Y$ be an object of $\mathcal{A}$ and $X$ be a filtered twisted complex over $\mathcal{A}$. Then there is a filtered quasi-isomorphism of $A_\infty$-modules $$\widetilde{\mathcal{Y}}_{T_YX} \to \Tau_Y\widetilde{\mathcal{Y}}_X.$$
\end{lem}
\subsubsection{Functoriality of filtered twisted complexes}
Let $\mathcal{A}$ and $\mathcal{B}$ be filtered $A_\infty$-categories. Recall that, for any $s\geq 0$, we defined the filtered $A_\infty$-category $\text{fun}^{\text{LD};s}(\mathcal{A}, \mathcal{B})$ of $A_\infty$-functors with linear deviation rate $s$ in \S\ref{a:func-LD}.\\
We introduce the following notations: for $N\geq 1$ and $s\geq0$ we write \begin{equation}\label{eq:FTwNQ}
	FTw^N\mathcal{Q}_s= \text{fun}^{\text{LD};s}(FTw^N\mathcal{A}, FTw^N\mathcal{B})
\end{equation} as well as $$FTw\mathcal{Q}_s= \text{fun}^{\text{LD};s}(FTw\mathcal{A}, FTw\mathcal{B})$$ Recall that $FTw^N\mathcal{A}$ is the full $A_\infty$-subcategory ()of the $A_\infty$-category $FTw\mathcal{A}$ of filtered twisted complexes) admitting only twisted complexes of length $\leq N$ as objects (see \S\ref{ap:ftc}). In this section we define filtered $A_\infty$-functors $$FTw^N\colon \text{fun}^{\text{LD};s}(\mathcal{A}, \mathcal{B})\to FTw^N\mathcal{Q}_{Ns}$$ for any $s\geq 0$ and $N\geq 1$. Note that in the target category, the deviation is $Ns$. Unfortunately, these functors do not extend to a functor $FTw$ because the deviation of functors in the image depends on the length of the twisted complexes it is applied to, as will be apparent from the discussion below. The best we can do is to define functors $$\overline{FTw}\colon \text{fun}^{\text{LD};s}(\mathcal{A}, \mathcal{B})\to \overline{FTw\mathcal{Q}_s}$$ where $\overline{FTw\mathcal{Q}_s}$ is the $A_\infty$-category of $A_\infty$-functors from $FTw\mathcal{A}$ to $FTw\mathcal{B}$ such that, for any $N\geq 1$, when restricted to the subcategory $FTw^N\mathcal{A}$ they have linear deviation rate $Ns$.\\

Let $s\geq 0$ and $(\mathcal{F},s)$ be an $A_\infty$-functor with linear deviation $s$. Then $\mathcal{F}$ immediately extends to an $A_\infty$-functor $\mathcal{F}^\Sigma\colon \mathcal{A}^\Sigma\to \mathcal{B}^\Sigma$ with deviation $s$ between the shift completions of $\mathcal{A}$ and $\mathcal{B}$. We define the extension $(\mathcal{F},s)^\oplus\colon \mathcal{A}^\oplus\to \mathcal{B}^\oplus$ to the filtered additive enlargments as follows (we drop the $s$ from the notation, but the induced functor depends on this choice):
\begin{enumerate}
	\item on an object $\overline{L}=\bigoplus_{i=1}^n\Sigma^{r_i}L_i[d_i]$, $\mathcal{F}^\oplus$ acts as \begin{equation}\label{eq:FTw}
		\mathcal{F}^\oplus\left(\bigoplus_{i=1}^n\Sigma^{r_i}L_i[d_i]\right):= \bigoplus_{i=1}^n\Sigma^{r_i-(i-1)s}\mathcal{F}(L_i)[d_i]
	\end{equation}
	\item given objects $\overline{L}_0,\ldots, \overline{L}_d$ of $\mathcal{A}^\oplus$ and morphisms $f^i\in \mathcal{A}^\oplus(\overline{L_{i-1}},\overline{L_i})$ for any $i=1,\ldots, d$, we set $\mathcal{F}^\oplus_d(f^1,\ldots,f^d)\in \mathcal{B}^\oplus(\mathcal{F}\overline{L}_0,\mathcal{F}\overline{L}_d)$ to be the matrix with $(i,j)$-entry equal to $$\sum_{i_1=1}^{|\overline{L}_1|}\cdots \sum_{i_{d-1}=1}^{|\overline{L}_{d-1}|}\mathcal{F}_d\left(f^1_{ii_1},\ldots, f^d_{i_{d-1}j}\right)$$ for $i=1,\ldots, |\overline{L_0}|$ and $j=1,\ldots, |\overline{L_d}|$.
\end{enumerate}
We can then easily extend $\mathcal{F}$ to an $A_\infty$-functor $FTw\mathcal{F}\colon FTw\mathcal{A}\to FTw\mathcal{B}$ by setting $$FTw\mathcal{F}(\overline{L}):= \mathcal{F}^\oplus(\overline{L})\text{  and  } q_{Tw\mathcal{F}(\overline{L})} := \sum_{d\geq 1}\mathcal{F}^\oplus_d(q_{\overline{L}},\ldots, q_{\overline{L}})$$ for any object $(\overline{L},q_{\overline{L}})$ of $FTw\mathcal{A}$, and, given objects $(\overline{L}_0,q_0),\ldots, (\overline{L}_d,q_d)$ of $FTw\mathcal{A}$ and morphisms $f^i\in FTw\mathcal{A}(\overline{L_{i-1}},\overline{L_i})$ for any $i=1,\ldots, d$, we set
$$(FTw\mathcal{F})_d(f^1,\ldots, f^d) := \sum_{k_0,\ldots, k_d\geq 0}\mathcal{F}^\oplus_{d+k_0+\cdots + k_d}\left(q_0^{\otimes k_0}, f^1,q_1^{\otimes k_1},\ldots, q_{d-1}^{\otimes k_{d-1}},f^d, q_d^{\otimes k_d}\right)$$
\begin{rem}
	\begin{enumerate}
		\item The sums in the definition of $q_{Tw\mathcal{F}(\overline{L})}$ and $FTw\mathcal{F}_d$ are finite. 
		\item The $s$-shifts in the image of $\overline{L}$ in (\ref{eq:FTw}) is essential in order for the image $q_{Tw\mathcal{F}(\overline{L})}$ of $q_{\overline{L}}$ to be at filtration level $\leq 0$, as we prove in the next lemma. Moreover, we remark here again that the induced functors $\mathcal{F}^\oplus$ and $FTw\mathcal{F}$ obviously depend on the choice of deviation $s$, although this fact is not explicitly reflected by the notation.
	\end{enumerate}
\end{rem}
\begin{lem}
	Let $(\mathcal{F},s)$ be an object of $\text{fun}^{\text{LD};s}(\mathcal{A}, \mathcal{B})$. Let $(\overline{L},q_{\overline{L}})$ be a filtered twisted complex. Then its image $(FTw\mathcal{F}(\overline{L}),q_{Tw\mathcal{F}(\overline{L})})$ under $FTw\mathcal{F}$ is a filtered twisted complex as well.
\end{lem}
\begin{rem}
	The fact that the image $FTw\mathcal{F}(\overline{L})$ preserves the structure of a filtered twisted complex although $\mathcal{F}$ is not in general a filtered $A_\infty$-functor may seem a bit counterintuitive, but might be taken as an hint that functors with linear deviation behave well with respect to filtered homological constructions and are hence the right setting.
\end{rem}
\begin{proof}
	It is easy to see that $q_{Tw\mathcal{F}(\overline{L})}$ satisfies the Maurer-Cartan equation. It remains to check that it lies at non-positive filtration level. We write $q$ for $q_{\overline{L}}$. Let $d\in \{1,\ldots, |\overline{L}|-1\}$ and consider the $(i,j)$-entry of $\mathcal{F}^\oplus_d(q,\ldots, q)$, which equals $$\sum_{i<i_1<\cdots < i_{d-1}<j} \mathcal{F}_d(q_{ii_1},\ldots, q_{i_{d-1}j}).$$ Assume that $\mathcal{F}_d(q_{ii_1},\ldots, q_{i_{d-1}j})$ is non-zero for some $i_1,\ldots, i_{d-1}$; in particular, $j-i\geq d$, as otherwise, by triangularity of $q$, this term would be zero. Note that $\mathcal{F}_d(q_{ii_1},\ldots, q_{i_{d-1}j})$ lies in $$(\mathcal{B}^\Sigma)^{\leq ds} (\Sigma^{r_i}\mathcal{F}(L_i)[d_i], \Sigma^{r_j}\mathcal{F}(L_j)[d_j])\subset (\mathcal{B}^\Sigma)^{\leq (j-i)s} (\Sigma^{r_i}\mathcal{F}(L_i)[d_i], \Sigma^{r_j}\mathcal{F}(L_j)[d_j])$$ which is isomorphic to $$ (\mathcal{B}^\Sigma)^{\leq 0} (\Sigma^{r_i-(i-1)s}\mathcal{F}(L_i)[d_i], \Sigma^{r_j-(j-1)s}\mathcal{F}(L_j)[d_j]).$$ By our definition of the action of $FTw\mathcal{F}$ on objects of $FTw\mathcal{A}$, it follows directly that $q_{Tw\mathcal{F}(\overline{L})}$ lies at vanishing filtration level.
\end{proof}
\begin{rem}
	Since the induced functor $FTw\mathcal{F}$ preserves lengths of twisted complexes by definition, it restricts to a functor $Tw^N\mathcal{A}\to Tw^N\mathcal{B}$. We denote this functor by $FTw^N\mathcal{F}$, i.e. highlighting the $N$, because of the following proposition.
\end{rem}
\begin{prop}
	Let $N\geq 1$ and let $(\mathcal{F},s)$ be an object of $\text{fun}^{\text{LD};s}(\mathcal{A}, \mathcal{B})$. Then the induced functor $$Tw^N\mathcal{F}\colon FTw^N\mathcal{A}\to FTw^N\mathcal{B}$$ is an $A_\infty$-functor with deviation rate $N\cdot s$, that is, an object of $FTw^NQ_{Ns}$ (see (\ref{eq:FTwNQ})).
\end{prop}
\begin{proof}
	Consider objects $(\overline{L^0},q_0),\ldots, (\overline{L^d},q_d)$ of $FTw^N\mathcal{A}$ of length $n^0,\ldots, n^d\leq N$ respectively, and morphisms $f^i\in FTw\mathcal{A}^{\alpha_i}(\overline{L_{i-1}},\overline{L_i})$ for any $i=1,\ldots,d$. The "most shifted" (in general) non-zero contribution to the matrix $$FTw\mathcal{F}_d(f^1,\ldots, f^d)\in FTw\mathcal{B}(Tw\mathcal{F}\overline{L_0},Tw\mathcal{F}\overline{L_d})$$ is the $(1,n^d)$ entry of the summand $$\mathcal{F}^\oplus_{d + (n^0-1)+\cdots + (n^d-1)}(q_0^{n^0-1},f^1,\ldots, f^d,q_d^{n^d-1}).$$
	This term lies at filtration level $$\leq \sum_{i=1}^d\alpha_i + \sum_{j=0}^d(n^j-1)s+ds = \sum_{i=1}^d\alpha_i + \sum_{j=0}^dn^js-s$$ in $\mathcal{B}^\Sigma(\Sigma^{r_1}\mathcal{F}(L^0_1),\Sigma^{r_{n^d}}\mathcal{F}(L^d_{n^d}))$, i.e. at level $$\leq \sum_{i=1}^d\alpha_i + \sum_{j=0}^{d-1}n^js$$ in 
	$\mathcal{B}^\Sigma(\Sigma^{r^0_1}\mathcal{F}(L^0_1),\Sigma^{r^d_{n^d}-(n^d-1)s}\mathcal{F}(L^d_{n^d}))$. As by assumption $n^0,\ldots, n^d\leq N$, we have $\sum_{j=0}^{d-1}n^js\leq dNs$ and the claim follows.
	
	%$$\mathcal{F}_{1 + (|\overline{L_0}|-1)+\cdots + (|\overline{L_d}|-1)}(q_0^{|\overline{L_0}|-1},f^1,\ldots, f^d,q_d^{|\overline{L_d}|-1})$$ Fix $i\in \{1,\ldots, |\overline{L_0}|\}$ and $j\in \{1,\ldots, |\overline{L_d}|\}$ and consider the $(i,j)$-entry of some summand $\mathcal{F}^\oplus_{d+k_0+\cdots + k_d}\left(q_0^{\otimes k_0}, f^1,q_1^{\otimes k_1},\ldots, q_{d-1}^{\otimes k_{d-1}},f^d, q_d^{\otimes k_d}\right)$ of $Tw\mathcal{F}_d(f^1,\ldots, f^d)$. 
\end{proof}
\begin{rem}
	The fact that $FTw\mathcal{F}$ acts on twisted complexes shifting summand by a quantity depending on the order is needed in order to have induced functors with shift that is sublinear on the order. The fact that the shift depend on the order seems counterintuitive, but it is in line with upper triangularity of differentials, which of course depends on the order of the summands.
\end{rem}

So far, we defined the $A_\infty$-functor $FTw$ just at the level of objects. We now define the first order term $FTw_1$. Let $(\mathcal{F},s)$ and $(\mathcal{G},s)$ be objects of $\text{fun}^{\text{LD};s}(\mathcal{A}, \mathcal{B})$ and consider an $A_\infty$-pre-natural transformation $T$ from $(\mathcal{F},s)$ to $(\mathcal{G},s)$ with shift $\rho$ (see \S\ref{a:func-LD} for the definition). We define the '$A_\infty$-pre-natural transformation'\footnote{The quotation marks are there as this is formally a natural transformation between functors that are not well-defined.} $FTw_1T$ with $d$th term, $d\geq 0$, $$(FTw_1T)_d(f_1,\ldots, f_d):= \sum_{k_0,\ldots, k_d\geq 0}T^\oplus_{d+k_0+\cdots + k_d}\left(q_0^{\otimes k_0}, f^1,q_1^{\otimes k_1},\ldots, q_{d-1}^{\otimes k_{d-1}},f^d, q_d^{\otimes k_d}\right)$$ where $T^\oplus$ is the matrixwise extension of $T$, defined analogously to the case of functors above. As in the case of functors, we can restrict to subcategories of twisted complexes of length $\leq N$ and define $FTw_1^NT$ for any $N\geq 1$. Taking filtrations into account, we have the following result.
\begin{lem}
	Let $(\mathcal{F},s)$ and $(\mathcal{G},s)$ be objects of $\text{fun}^{\text{LD};s}(\mathcal{A}, \mathcal{B})$ and consider an $A_\infty$-pre-natural transformation $T\colon (\mathcal{F},s)\to(\mathcal{G},s)$ of shift $\rho$ between them. Let $N\geq 1$. Then $FTw^N_1T$ is an $A_\infty$-pre-natural transformation between $FTw^N\mathcal{F}$ and $FTw^N\mathcal{G}$ of shift $\leq \rho$, that is $$FTw^N_1T\in FTw^NQ_{Ns}(FTw^N\mathcal{F},FTw^N\mathcal{G})^\rho.$$
\end{lem}
\begin{rem}
	In other words, the above lemma tells us that $FTw_1$ sends $A_\infty$-pre-natural transformations with linear deviation $s$ and shift $\leq \rho$ to $A_\infty$-pre-natural transformations with linear deviation $Ns$ and shift $\rho$.
\end{rem}
%The same is true if we consider a natural transformation between objects $(\mathcal{F},s)$ and $(\mathcal{G},r)$ of $\mathcal{Q}$. \noindent This generalizes to:
We define the higher order terms $FTw_d$, $d\geq 2$, to be zero. In particular, we proved the following result, which was announced at the beginning of this section.
\begin{cor}
	For any $N\geq 1$ and $s\geq 0$, the $A_\infty$-functor $$FTw^N\colon \text{fun}^{\text{LD};s}(\mathcal{A}, \mathcal{B})\to \text{fun}^{\text{LD};Ns}(FTw^N\mathcal{A},FTw^N\mathcal{B}) = FTw^N\mathcal{Q}_{Ns}$$ is filtered.
\end{cor} 

